\chapter{Trajectory Prediction: From Constant Velocity to LSTMs and Transformers}
\label{ch:ch20}
% Function names for pseudocode are prefixed with "Pred" to stay unique
% across the book.
\SetKwFunction{PredCV}{ConstantVelocity}
\SetKwFunction{PredCA}{ConstantAcceleration}
\SetKwFunction{PredKF}{KalmanExtrapolate}
\SetKwFunction{PredCell}{LSTMCell}
\SetKwFunction{PredPredict}{Seq2SeqPredict}
\SetKwFunction{PredTrain}{TrainSeq2Seq}
\SetKwFunction{PredEvaluate}{EvaluatePredictor}
\SetKwFunction{PredAttention}{Attention}

\begin{objectives}
\begin{itemize}
  \item State the trajectory-prediction problem (observe $T_{\mathrm{obs}}$ positions, predict the next $T_{\mathrm{pred}}$), define ADE, FDE, $\mathrm{minADE}_K$ and collision-based metrics exactly, and design an evaluation that cannot fool you.
  \item Implement and tune the baselines that a learned model must beat: constant velocity, constant acceleration and Kalman-filter extrapolation, and explain why constant velocity is so hard to beat at short horizons.
  \item Write down the LSTM cell with its three gates, trace one step by hand, and build a sequence-to-sequence predictor with an agent-centred input frame, a step-by-step decoder, a Gaussian output and a training loop with Adam and early stopping.
  \item Explain scaled dot-product attention, multi-head attention and positional encodings, and how attention across agents makes a predictor interaction-aware.
  \item Turn a predicted distribution into a time-indexed occupancy region and feed it to the local, replanning and control layers of the hybrid system.
\end{itemize}
\end{objectives}

% ---------------------------------------------------------------------
\section{Why this matters for a drone swarm}
\label{sec:ch20-motivation}

The tracker of \cref{ch:ch18} hands the planner a clean estimate of where an
unknown drone \emph{is}. That is not enough. A drone that flies at $3$~m/s
covers $7.5$~m in the $2.5$~s that a swarm member needs to brake, turn and
re-establish its formation, and the collision that matters is the one at the
end of those $2.5$~s, not the one now. The local safety layer of
\cref{ch:ch24} therefore asks a different question: where will the intruder
be at $t + 0.2$~s, $t + 0.4$~s, \dots, $t + 2.4$~s, and how sure are we?
This is \emph{trajectory prediction}, the second half of the prediction
layer of the hybrid architecture.

The simplest answer, ``where it is now plus its current velocity times the
elapsed time'', is the constant-velocity extrapolation that the Kalman
filter of \cref{ch:ch18} already produces when it runs its prediction step
several times in a row. For a drone in steady flight this is exact, and
for short horizons it is remarkably hard to beat: a large survey of human
motion prediction~\cite{rudenko2020human} and a widely discussed follow-up
study~\cite{schoeller2020constant} both found that a carefully tuned
constant-velocity model matches or beats many published neural predictors
on the standard pedestrian benchmarks. The reason is not that learning is
useless but that most of a trajectory \emph{is} straight flight, and the
learned models were often compared against a badly tuned baseline. The
interesting cases are the others: the intruder is banking into a turn, it
has just begun an evasive manoeuvre, or it is one of several drones whose
paths interact. There the history of the last second reveals what happens
next, and a model that has learned the typical shapes of manoeuvres from
data predicts the turn while the straight-line model flies off along the
tangent.

The two ingredients of the Week~10 milestone of the study plan
(\cref{ch:appA}) are therefore a strong baseline and an honest measurement:
you should be able to \emph{quantify} whether learning improves prediction,
per horizon and per kind of manoeuvre, before you let a learned predictor
decide the safety margins of a swarm. The chapter builds both. It also
prepares what \cref{ch:ch24} needs: not a single predicted path but a
distribution over paths, turned into a growing ellipse per time step, so
that \orca (\cref{ch:ch13}) can inflate the intruder's radius, \dstarlite
(\cref{ch:ch05}) can raise the cost of the cells it may occupy, and the MPC
of \cref{ch:ch21} can impose a chance constraint.

The chapter proceeds as follows. \Cref{sec:ch20-intuition} gives the idea and
\cref{sec:ch20-problem} the problem, the metrics and the evaluation
protocol. \Cref{sec:ch20-baselines} presents the physics baselines and
\cref{sec:ch20-lstm} the LSTM cell, the sequence-to-sequence predictor and
the Social-LSTM idea. \Cref{sec:ch20-transformer} covers attention and
Transformers, \cref{sec:ch20-multimodal} multimodality and uncertainty.
\Cref{sec:ch20-example} is the experiment of the chapter, run with the
book's own NumPy code, and \cref{sec:ch20-properties} proves the few results
that make the experiment interpretable: when constant velocity is optimal,
how noise is amplified, why a squared-error predictor averages the modes,
and what the sequence models cost. \Cref{sec:ch20-variants},
\cref{sec:ch20-implementation} and \cref{sec:ch20-drone} close with variants,
implementation notes and the place of prediction in the drone system.

% ---------------------------------------------------------------------
\section{The idea in plain words}
\label{sec:ch20-intuition}

\Cref{fig:ch20-idea} shows the situation. A drone has been observed at eight
time steps; the black dots are its measured positions, $0.2$~s apart. We
want the next twelve positions. Three things are visible in the picture and
they organise the whole chapter.

\begin{figure}[tb]
  \centering
  \inputfigure{ch20/idea}
  \caption{The prediction problem. The last $T_{\mathrm{obs}} = 8$ positions of
  an intruder are observed; the next $T_{\mathrm{pred}} = 12$ must be predicted.
  The blue line is a prediction, the green line the future that actually
  happens, and the red segments are the displacement errors $e_h$ whose
  average is the ADE and whose last value is the FDE. The ellipses say how
  uncertain the prediction is at each step; they grow with the horizon. The
  dashed orange line is a second future that would have been just as
  plausible given the history: predictions are distributions, not lines.}
  \label{fig:ch20-idea}
\end{figure}

First, the past constrains the future. The drone will not jump; its next
position is near its last one, and its next velocity is near its last
velocity. A predictor that uses only this fact is the constant-velocity
model, and for the first few steps nothing does much better.

Second, the past reveals the manoeuvre. If the last four displacements
rotate steadily to the left, the drone is turning, and it will probably
keep turning for a while. If the rotation started suddenly two steps ago,
it may be an evasive manoeuvre, which typically lasts a second and ends in
straight flight. Which of these patterns occur, how long they last and how
they end is a property of the \emph{population} of trajectories, and it can
be learned from recorded data. A recurrent network that reads the history
one step at a time, keeps a memory of what it has seen and writes the
future one step at a time is the natural tool; a Transformer that lets
every time step and every neighbouring agent attend to every other is the
modern alternative.

Third, the future is uncertain, and the uncertainty grows. Part of it is
noise; part of it is a genuine choice the intruder has not made yet. At a
junction it may go left or right, and no model can know which. A predictor
that outputs a single line is then wrong twice: its line lies between the
branches, and it does not say how unsure it is. The planner needs the
spread, in the form of an ellipse per step or a handful of sampled futures.

\begin{keyidea}
Predict the next $T_{\mathrm{pred}}$ positions from the last
$T_{\mathrm{obs}}$, in the intruder's own frame, and measure the error at
every horizon. Constant velocity is the baseline that must be beaten;
learned sequence models beat it where the history reveals a manoeuvre;
nothing beats it on straight flight or on a manoeuvre that has not started
yet. Always predict a distribution, because the planner needs the spread as
much as the mean.
\end{keyidea}

% ---------------------------------------------------------------------
\section{Problem statement, metrics and evaluation protocol}
\label{sec:ch20-problem}

\subsection{The prediction problem}

\begin{definition}[Trajectory prediction]
\label{def:ch20-problem}
A \textbf{trajectory}\index{trajectory prediction} is a sequence of positions
$\pos_t \in \R^2$ (or $\R^3$) sampled every $\dt$ seconds. At the current time
$T$ the \textbf{observation window} of length $T_{\mathrm{obs}}$ contains the
positions $\pos_{T - T_{\mathrm{obs}} + 1}, \dots, \pos_T$ (in practice noisy
measurements of them). A \textbf{predictor} maps the observation window,
and possibly the windows of other agents and a map, to predicted positions
$\hat{\pos}_{T+1}, \dots, \hat{\pos}_{T+T_{\mathrm{pred}}}$ over the
\textbf{prediction horizon}\index{prediction horizon} $T_{\mathrm{pred}}$. The
integer $h \in \set{1, \dots, T_{\mathrm{pred}}}$ is the \textbf{horizon
step}; $h\,\dt$ is the look-ahead time.
\end{definition}

Three distinctions matter. A \textbf{single-agent} predictor sees one
history at a time; an \textbf{interaction-aware}\index{interaction-aware
prediction} predictor also sees the neighbours, because drones and people
avoid each other. A \textbf{deterministic} predictor outputs one trajectory;
a \textbf{probabilistic} one outputs a distribution, either as a mean and a
covariance per step or as $K$ sampled trajectories, and a
\textbf{multimodal}\index{multimodal prediction} one can put mass on several
distinct futures. Throughout the chapter $\dt = 0.2$~s, $T_{\mathrm{obs}} = 8$
($1.6$~s) and $T_{\mathrm{pred}} = 12$ ($2.4$~s), which mirrors the
$8/12$-step convention of the pedestrian benchmarks
ETH~\cite{pellegrini2009you} and UCY~\cite{lerner2007crowds} at a faster
sampling rate suited to drones.

\subsection{Metrics}

\begin{definition}[Displacement errors]
\label{def:ch20-metrics}
For one trajectory with true future $\pos_{T+h}$ and prediction
$\hat{\pos}_{T+h}$, the \textbf{displacement error} at step $h$ is
$e_h = \norm{\hat{\pos}_{T+h} - \pos_{T+h}}$. The \textbf{average displacement
error}\index{ADE (average displacement error)} and the \textbf{final displacement
error}\index{FDE (final displacement error)} are
\begin{equation}
  \mathrm{ADE} = \frac{1}{T_{\mathrm{pred}}} \sum_{h=1}^{T_{\mathrm{pred}}} e_h,
  \qquad
  \mathrm{FDE} = e_{T_{\mathrm{pred}}} .
  \label{eq:ch20-ade-fde}
\end{equation}
Over a test set the ADE and FDE are the means over its trajectories. The
\textbf{per-horizon} versions are $\mathrm{ADE}_h = \frac{1}{h}\sum_{j \le h}
e_j$ and $\mathrm{FDE}_h = e_h$, again averaged over the test set. For a
predictor that outputs $K$ sampled trajectories with errors
$e_h^{(1)}, \dots, e_h^{(K)}$,
\begin{equation}
  \mathrm{minADE}_K = \min_{k} \frac{1}{T_{\mathrm{pred}}} \sum_{h} e_h^{(k)},
  \qquad
  \mathrm{minFDE}_K = \min_{k} e_{T_{\mathrm{pred}}}^{(k)},
  \label{eq:ch20-minade}
\end{equation}
\index{minADE@$\mathrm{minADE}_K$}
each averaged over the test set. The \textbf{miss rate} is the fraction of
trajectories with $\mathrm{FDE} > \tau_{\mathrm{miss}}$ for a threshold such
as $\tau_{\mathrm{miss}} = 1$~m.
\end{definition}

The ADE rewards a prediction that is close on average; the FDE isolates the
end of the horizon, where errors are largest and where the planner's safety
margin is decided. $\mathrm{minADE}_K$ is generous by construction: it scores
a sampler by its \emph{best} of $K$ guesses, which is the right measure of
whether the true future is covered by the samples but says nothing about
how the planner should choose among them. A paper that reports
$\mathrm{minADE}_{20}$ next to the ADE of a deterministic baseline compares
two different things.

Displacement errors ignore physics. A prediction that passes through a wall
or through another drone has a small ADE if it ends near the truth, but a
planner that trusts it will be surprised. \textbf{Collision-based
metrics}\index{collision rate!of predictions} count such predictions: the
fraction of predicted trajectories that come closer than $d_{\min}$ to a
static obstacle, or to the true trajectory of another agent at the same
time step. The book's evaluation chapter, \cref{ch:ch25}, measures the
collision rate of the \emph{planner} that consumes the predictions, which is
the metric that finally matters.

\subsection{An honest evaluation protocol}
\label{sec:ch20-protocol}

Prediction results are easy to inflate, mostly by accident.
\Cref{alg:ch20-evaluate} fixes the protocol used in this chapter.

\begin{algorithm}[htb]
\caption{Evaluation protocol for a trajectory predictor}
\label{alg:ch20-evaluate}
\KwIn{recorded trajectories grouped by scene, candidate predictors, baselines with their tuning parameters, $T_{\mathrm{obs}}$, $T_{\mathrm{pred}}$, $\dt$}
\KwOut{per-horizon $\mathrm{ADE}_h$, $\mathrm{FDE}_h$ for every predictor and every subset}
\Fn{\PredEvaluate{scenes}}{
  split the \emph{scenes} into train, validation and test sets\label{alg:ch20-evaluate:split}\tcp*{never split windows of one trajectory}
  cut every trajectory into windows of $T_{\mathrm{obs}} + T_{\mathrm{pred}}$ samples at the common $\dt$\;
  compute normalisation statistics on the training set only\label{alg:ch20-evaluate:norm}\;
  \ForEach{baseline}{tune its parameters ($k$, $q$, \dots) on the validation set\label{alg:ch20-evaluate:tune}\;}
  \ForEach{learned predictor}{train on the training set; select epochs and hyperparameters on the validation set\;}
  \ForEach{predictor, subset of the test set (manoeuvre type, density, \dots)}{
    report $\mathrm{ADE}_h$ and $\mathrm{FDE}_h$ for $h = 1, \dots, T_{\mathrm{pred}}$, the miss rate and, for samplers, $\mathrm{minADE}_K$ with $K$ stated\label{alg:ch20-evaluate:report}\;
  }
  repeat with several seeds and report the spread\;
}
\end{algorithm}

Line~\ref{alg:ch20-evaluate:split} is the one most often violated. A
recording is usually cut into overlapping windows; if windows are shuffled
and split at random, the test set contains windows whose observed part was
the predicted part of a training window, and the model has seen the
answers. This is \textbf{data leakage}\index{data leakage}: split by scene
or, for a single long recording, by time, and keep the scenes that contain
the manoeuvres you care about in the test set. Line~\ref{alg:ch20-evaluate:norm}
prevents a subtler leak through the normalisation. Line~\ref{alg:ch20-evaluate:tune}
gives the baseline the same care as the learned model: the velocity
estimate of the constant-velocity model has a window length, the Kalman
filter has a process noise, and both change the ADE by a large factor
(\cref{sec:ch20-example}). Line~\ref{alg:ch20-evaluate:report} asks for
curves rather than numbers: an ADE at a single, unstated horizon hides
whether a model is good at $0.4$~s or at $2.4$~s, and per-subset results
reveal where the model wins and where it does not.

% ---------------------------------------------------------------------
\section{The baselines: constant velocity, constant acceleration and Kalman extrapolation}
\label{sec:ch20-baselines}

\begin{algorithm}[htb]
\caption{Physics baselines for a horizon of $H$ steps}
\label{alg:ch20-baselines}
\KwIn{observed positions $\meas_1, \dots, \meas_{T_{\mathrm{obs}}}$ (noisy), step $\dt$, horizon $H$, window $k$, filter matrices $\mat{F}, \mat{Q}, \mat{H}, \mat{R}$}
\KwOut{predicted positions $\hat{\pos}_1, \dots, \hat{\pos}_H$ (and covariances $\mat{\Sigma}_1, \dots, \mat{\Sigma}_H$ for the filter)}
\Fn{\PredCV{$\meas$, $H$, $k$}}{
  $\hat{\vel} \gets (\meas_{T_{\mathrm{obs}}} - \meas_{T_{\mathrm{obs}} - k}) / (k\,\dt)$\label{alg:ch20-baselines:vel}\tcp*{mean velocity over the last $k$ intervals}
  \KwRet $\hat{\pos}_h = \meas_{T_{\mathrm{obs}}} + h\,\dt\,\hat{\vel}$ for $h = 1, \dots, H$\;
}
\Fn{\PredCA{$\meas$, $H$, $k$}}{
  $\vel_1 \gets (\meas_{T_{\mathrm{obs}}} - \meas_{T_{\mathrm{obs}} - k}) / (k\,\dt)$; \quad $\vel_0 \gets (\meas_{T_{\mathrm{obs}} - k} - \meas_{T_{\mathrm{obs}} - 2k}) / (k\,\dt)$\;
  $\hat{\acc} \gets (\vel_1 - \vel_0)/(k\,\dt)$; \quad $\hat{\vel} \gets \vel_1 + \tfrac{1}{2}\hat{\acc}\,k\,\dt$\label{alg:ch20-baselines:acc}\tcp*{$\vel_1$ is the velocity at the midpoint}
  \KwRet $\hat{\pos}_h = \meas_{T_{\mathrm{obs}}} + h\,\dt\,\hat{\vel} + \tfrac{1}{2}(h\,\dt)^2\hat{\acc}$ for $h = 1, \dots, H$\;
}
\Fn{\PredKF{$\meas$, $H$, $\mat{F}$, $\mat{Q}$, $\mat{H}$, $\mat{R}$}}{
  initialise $\hat{\state} \gets (\meas_2,\, (\meas_2 - \meas_1)/\dt)$ and $\mat{P}$ from $\mat{R}$\;
  \For{$t = 3, \dots, T_{\mathrm{obs}}$}{run one predict and one update step of the Kalman filter of \cref{ch:ch18} with measurement $\meas_t$\label{alg:ch20-baselines:filter}\;}
  \For{$h = 1, \dots, H$}{
    $\hat{\state} \gets \mat{F}\hat{\state}$; \quad $\mat{P} \gets \mat{F}\mat{P}\mat{F}\T + \mat{Q}$\label{alg:ch20-baselines:extrapolate}\tcp*{predict step only: no measurement}
    $\hat{\pos}_h \gets \mat{H}\hat{\state}$; \quad $\mat{\Sigma}_h \gets \mat{H}\mat{P}\mat{H}\T$\;
  }
  \KwRet $(\hat{\pos}_h, \mat{\Sigma}_h)_{h=1}^{H}$\;
}
\end{algorithm}

\Cref{alg:ch20-baselines} lists the three predictors that every learned model
must beat, and \cref{fig:ch20-baselines} shows them on a gently turning test
trajectory of the experiment in \cref{sec:ch20-example}.

\paragraph{Constant velocity (CV).}
\index{constant-velocity model!as prediction baseline}
The velocity is estimated from the last $k$ intervals
(line~\ref{alg:ch20-baselines:vel}) and held constant. With $k = 1$ the
estimate uses the last two measurements only and is very noisy: each
measurement error of standard deviation $\sigma$ enters the velocity
divided by $\dt$, and \cref{thm:ch20-cv-noise} shows that the position
error at horizon $h$ then grows like $h\sigma/\dt$. A larger $k$ averages
the noise but lags behind a turn. The window $k$ is the tuning parameter of
the baseline and must be chosen on the validation set; in the experiment
below the naive $k = 1$ is $88\,\%$ worse than the tuned $k = 2$ on straight
flight. This is the \emph{poorly tuned baseline} that many papers beat.

\paragraph{Constant acceleration (CA).}
\index{constant-acceleration model}
Two successive velocity estimates give an acceleration
(line~\ref{alg:ch20-baselines:acc}). On a circular turn the acceleration
points to the centre and rotates with the drone, so extrapolating it as a
constant is right for one or two steps and then overshoots: the CA
prediction in \cref{fig:ch20-baselines} cuts inside the turn and ends $4$~m
from the truth, further than either CV prediction. The second finite
difference also triples the noise amplification. CA is included because it
is the natural next term and because its failure shows that
``more physics'' in the extrapolation is not automatically better.

\paragraph{Kalman extrapolation (KF).}
\index{Kalman filter!extrapolation}\index{Kalman extrapolation}
The constant-velocity Kalman filter of \cref{ch:ch18}, with state
$\state = (\pos, \vel) \in \R^4$, transition
$\mat{F} = \begin{psmallmatrix} \mat{I} & \dt\,\mat{I} \\ \mat{0} & \mat{I}
\end{psmallmatrix}$, measurement $\mat{H} = (\mat{I}\;\;\mat{0})$, measurement
noise $\mat{R} = \sigma^2\mat{I}$ and the white-noise-acceleration process
noise
\begin{equation}
  \mat{Q} = q \begin{pmatrix} \tfrac{1}{3}\dt^3\,\mat{I} & \tfrac{1}{2}\dt^2\,\mat{I} \\[2pt]
                              \tfrac{1}{2}\dt^2\,\mat{I} & \dt\,\mat{I} \end{pmatrix},
  \label{eq:ch20-q}
\end{equation}
is run over the observed window and then \emph{extrapolated}: the predict
step is repeated $H$ times without an update
(line~\ref{alg:ch20-baselines:extrapolate}). The mean is a constant-velocity
extrapolation of the filtered state, so the filter is a CV model whose
velocity estimate weighs the whole window optimally under its noise model;
the covariance $\mat{\Sigma}_h$ grows with $h$ and is the first prediction in
this chapter that comes with an uncertainty
(\cref{fig:ch20-baselines}). The process-noise intensity $q$ is its tuning
parameter: small $q$ trusts the straight-line model and smooths the noise,
large $q$ follows the last measurements and reacts to turns. Because
$\mat{F}$, $\mat{Q}$, $\mat{H}$ and $\mat{R}$ do not depend on the data, the
covariance $\mat{P}$ and the gain are the same for every trajectory, which
makes the extrapolation of a whole test set a handful of matrix products.

\begin{figure}[tb]
  \centering
  \inputfigure{ch20/baselines}
  \caption{The physics baselines on a gently turning test trajectory
  (turn rate $0.17$~rad/s, speed $2.9$~m/s, the turn with the median error of
  the learned model). Both CV predictions leave the turn along a tangent;
  the constant-acceleration prediction (dashed) extrapolates the rotating
  centripetal acceleration as a constant and overshoots, ending $4.0$~m from
  the truth against $1.5$~m for the tuned CV. The Kalman extrapolation is a
  smoothed CV prediction with a $95\,\%$ ellipse at $h = 4, 8, 12$ that grows
  with the horizon. On this one trajectory the naive $k = 1$ happens to beat
  the tuned $k = 2$ ($0.8$ against $1.5$~m at $h = 12$); single examples prove
  nothing, which is why \cref{tab:ch20-results} averages over $800$.}
  \label{fig:ch20-baselines}
\end{figure}

None of the three knows that turns end, that evasive manoeuvres have a
typical duration, or that two drones on a collision course will both
deviate. That knowledge has to come from data.

% ---------------------------------------------------------------------
\section{Sequence models: the LSTM and sequence-to-sequence prediction}
\label{sec:ch20-lstm}

\subsection{The LSTM cell}
\label{sec:ch20-cell}

A \textbf{recurrent neural network}\index{recurrent neural network} reads a
sequence one element at a time and keeps a state vector that summarises
what it has read. The simplest form updates a hidden state $\vect{h}_t \in
\R^D$ as $\vect{h}_t = \tanh(\mat{W}[\vect{x}_t; \vect{h}_{t-1}] + \vect{b})$,
where $[\vect{x}_t; \vect{h}_{t-1}]$ is the concatenation of the current
input and the previous state. Training it means propagating the loss
gradient backwards through every step (\emph{backpropagation through
time}\index{backpropagation through time}, BPTT~\cite{werbos1990bptt}), and
each step multiplies the gradient by the Jacobian of the update, a matrix
whose norm is typically far from~$1$. After twenty steps the product has
either vanished or exploded, so the network cannot learn that something it
saw twenty steps ago matters now. Hochreiter and
Schmidhuber~\cite{hochreiter1997lstm} solved this with the \textbf{long
short-term memory}\index{LSTM (long short-term memory)} (LSTM) cell: a
second state, the \textbf{cell state}\index{LSTM (long short-term memory)!cell state}
$\vect{c}_t$, that is changed only by \emph{gated} additions, so that the
gradient can flow through it unchanged. The forget gate, which we use
throughout, was added by Gers, Schmidhuber and
Cummins~\cite{gers2000learning}.

\begin{figure}[tb]
  \centering
  \inputfigure{ch20/lstm-cell}
  \caption{One step of the LSTM cell. The cell state $\vect{c}$ runs along the
  top and is changed only by a multiplication with the forget gate and an
  addition of the gated candidate; the hidden state $\vect{h}$ is a gated,
  squashed read-out of it. All four gate vectors are computed from the same
  concatenated input $[\vect{x}_t; \vect{h}_{t-1}]$ with their own weights.}
  \label{fig:ch20-lstm-cell}
\end{figure}

\begin{definition}[LSTM cell]
\label{def:ch20-lstm}
An LSTM cell with input size $M$ and $D$ hidden units has weights
$\mat{W}_f, \mat{W}_i, \mat{W}_o, \mat{W}_c \in \R^{D \times (M + D)}$ and
biases $\vect{b}_f, \vect{b}_i, \vect{b}_o, \vect{b}_c \in \R^D$. Given the
input $\vect{x}_t \in \R^M$ and the previous states $\vect{h}_{t-1}, \vect{c}_{t-1}
\in \R^D$, one step computes, with $\vect{z}_t = [\vect{x}_t; \vect{h}_{t-1}]$,
the logistic function $\sigma(a) = 1/(1 + e^{-a})$ and $\odot$ the
element-wise product,
\begin{align}
  \vect{f}_t &= \sigma(\mat{W}_f \vect{z}_t + \vect{b}_f) && \text{forget gate,}\index{LSTM (long short-term memory)!gates} \label{eq:ch20-forget}\\
  \vect{i}_t &= \sigma(\mat{W}_i \vect{z}_t + \vect{b}_i) && \text{input gate,} \label{eq:ch20-input}\\
  \vect{o}_t &= \sigma(\mat{W}_o \vect{z}_t + \vect{b}_o) && \text{output gate,} \label{eq:ch20-output}\\
  \tilde{\vect{c}}_t &= \tanh(\mat{W}_c \vect{z}_t + \vect{b}_c) && \text{candidate values,} \label{eq:ch20-candidate}\\
  \vect{c}_t &= \vect{f}_t \odot \vect{c}_{t-1} + \vect{i}_t \odot \tilde{\vect{c}}_t && \text{cell state,} \label{eq:ch20-cell}\\
  \vect{h}_t &= \vect{o}_t \odot \tanh(\vect{c}_t) && \text{hidden state.} \label{eq:ch20-hidden}
\end{align}
\end{definition}

The two nonlinearities have different jobs. The logistic function maps
every pre-activation into $(0, 1)$, so the gate vectors are soft switches:
a forget-gate entry near $1$ keeps that component of the memory, near $0$
erases it; an input-gate entry decides how much of the candidate is
written; an output-gate entry decides how much of the memory is exposed to
the rest of the network. The hyperbolic tangent maps into $(-1, 1)$ and
produces the \emph{values}: the candidate that may be written and the
squashed memory that may be read. Because \cref{eq:ch20-cell} is a gated
sum rather than a squashed product, the derivative of $\vect{c}_t$ with
respect to $\vect{c}_{t-1}$ along the direct path is exactly
$\diag(\vect{f}_t)$ (\cref{thm:ch20-carousel}); with forget gates near $1$
the gradient neither vanishes nor explodes over long sequences, which is
why gating helps. \Cref{alg:ch20-lstm-cell} is the step as pseudocode; the
listing in \cref{sec:ch20-implementation} is its NumPy form together with
the backward pass.

\begin{algorithm}[htb]
\caption{One step of the LSTM cell}
\label{alg:ch20-lstm-cell}
\KwIn{input $\vect{x}_t$, previous states $\vect{h}_{t-1}$, $\vect{c}_{t-1}$, weights $\mat{W}$ (rows of $\mat{W}_f, \mat{W}_i, \mat{W}_o, \mat{W}_c$ stacked) and $\vect{b}$}
\KwOut{new states $\vect{h}_t$, $\vect{c}_t$}
\Fn{\PredCell{$\vect{x}_t$, $\vect{h}_{t-1}$, $\vect{c}_{t-1}$}}{
  $\vect{z} \gets [\vect{x}_t; \vect{h}_{t-1}]$; \quad $\vect{a} \gets \mat{W}\vect{z} + \vect{b}$\label{alg:ch20-lstm-cell:pre}\tcp*{$4D$ pre-activations at once}
  $\vect{f} \gets \sigma(\vect{a}_{1:D})$; $\vect{i} \gets \sigma(\vect{a}_{D+1:2D})$; $\vect{o} \gets \sigma(\vect{a}_{2D+1:3D})$; $\tilde{\vect{c}} \gets \tanh(\vect{a}_{3D+1:4D})$\label{alg:ch20-lstm-cell:gates}\;
  $\vect{c}_t \gets \vect{f} \odot \vect{c}_{t-1} + \vect{i} \odot \tilde{\vect{c}}$\label{alg:ch20-lstm-cell:cell}\;
  $\vect{h}_t \gets \vect{o} \odot \tanh(\vect{c}_t)$\;
  \KwRet $(\vect{h}_t, \vect{c}_t)$\;
}
\end{algorithm}

\begin{example}[One LSTM step by hand]
\label{ex:ch20-cell}
Take a cell with $M = 2$ inputs and $D = 2$ hidden units, the weights
\[
  \mat{W}_f = \begin{pmatrix} 0.5 & -0.5 & 0.3 & 0.1 \\ 0.2 & 0.4 & -0.3 & 0.6 \end{pmatrix},\;
  \mat{W}_i = \begin{pmatrix} 1.0 & 0.0 & 0.2 & -0.4 \\ -0.5 & 0.8 & 0.1 & 0.3 \end{pmatrix},\;
  \mat{W}_o = \begin{pmatrix} 0.3 & 0.3 & -0.6 & 0.2 \\ 0.7 & -0.2 & 0.4 & 0.5 \end{pmatrix},\;
  \mat{W}_c = \begin{pmatrix} 0.9 & -0.6 & 0.0 & 0.8 \\ -0.4 & 0.5 & 0.7 & -0.1 \end{pmatrix},
\]
$\vect{b}_f = (1, 1)$ (the usual initialisation that starts by remembering)
and all other biases zero. The input is $\vect{x}_t = (1.0, -0.5)$ and the
previous states are $\vect{h}_{t-1} = (0.2, -0.1)$ and $\vect{c}_{t-1} = (0.5,
-0.3)$, so $\vect{z}_t = (1.0, -0.5, 0.2, -0.1)$. \Cref{tab:ch20-cell} traces
\cref{alg:ch20-lstm-cell}; every number is produced by
\code{lstm\_cell\_example()} in \code{ch20\_prediction.py}. For the first
unit, the forget pre-activation is $0.5 \cdot 1.0 + (-0.5)(-0.5) + 0.3 \cdot
0.2 + 0.1 \cdot (-0.1) + 1 = 1.80$, so $f_1 = \sigma(1.80) = 0.8581$: the
unit keeps $86\,\%$ of its old memory $0.5$ and adds $i_1 \tilde{c}_1 = 0.7465
\cdot 0.8076 = 0.6029$, giving $c_1 = 0.4291 + 0.6029 = 1.0319$. The output
gate lets half of $\tanh(1.0319) = 0.7747$ through, $h_1 = 0.3893$. The
second unit has a small input gate ($0.2870$) and a negative candidate, so
its memory moves only from $-0.3$ to $-0.3447$.
\end{example}

\begin{table}[tb]
  \centering\small
  \caption{Trace of \cref{alg:ch20-lstm-cell} on \cref{ex:ch20-cell}: the
  pre-activations of line~\ref{alg:ch20-lstm-cell:pre}, the gate values of
  line~\ref{alg:ch20-lstm-cell:gates} (all in $(0,1)$, the candidate in
  $(-1,1)$) and the new states.}
  \label{tab:ch20-cell}
  \begin{tabular}{@{}lrrrrrr@{}}
  \toprule
  & \multicolumn{2}{c}{pre-activation $\vect{a}$} & \multicolumn{2}{c}{activation} & \multicolumn{2}{c}{state}\\
  \cmidrule(lr){2-3}\cmidrule(lr){4-5}\cmidrule(lr){6-7}
  & unit 1 & unit 2 & unit 1 & unit 2 & unit 1 & unit 2\\
  \midrule
  forget gate $\vect{f}_t$      & $1.80$ & $0.88$ & $0.8581$ & $0.7068$ & &\\
  input gate $\vect{i}_t$       & $1.08$ & $-0.91$ & $0.7465$ & $0.2870$ & &\\
  output gate $\vect{o}_t$      & $0.01$ & $0.83$ & $0.5025$ & $0.6964$ & &\\
  candidate $\tilde{\vect{c}}_t$ & $1.12$ & $-0.50$ & $0.8076$ & $-0.4621$ & &\\
  cell state $\vect{c}_t = \vect{f}_t \odot \vect{c}_{t-1} + \vect{i}_t \odot \tilde{\vect{c}}_t$ & & & & & $1.0319$ & $-0.3447$\\
  $\tanh(\vect{c}_t)$            & & & & & $0.7747$ & $-0.3316$\\
  hidden state $\vect{h}_t = \vect{o}_t \odot \tanh(\vect{c}_t)$ & & & & & $0.3893$ & $-0.2309$\\
  \bottomrule
  \end{tabular}
\end{table}

\subsection{A sequence-to-sequence predictor}
\label{sec:ch20-seq2seq}

\Cref{fig:ch20-seq2seq} shows how LSTM cells become a trajectory predictor.
An \textbf{encoder}\index{sequence-to-sequence model} LSTM reads the observed
history one step at a time; its final states $(\vect{h}, \vect{c})$ summarise
the history. A \textbf{decoder} LSTM starts from those states and rolls the
future out one step at a time: at every step a linear read-out layer turns
its hidden state into the next displacement, and that displacement is fed
back as the next input. Four design decisions make this work.

\begin{figure}[tb]
  \centering
  \inputfigure{ch20/seq2seq}
  \caption{The sequence-to-sequence predictor unrolled in time. The
  encoder cells (blue) read the seven observed displacements in the agent
  frame; their final state is handed to the decoder cells (orange), which
  produce, through the read-out $\mat{W}_{\mathrm{out}}$, the mean and standard
  deviation of the next displacement. Each decoder step is fed the mean of
  the previous one (dashed); teacher forcing replaces this feedback by the
  true displacement during training.}
  \label{fig:ch20-seq2seq}
\end{figure}

\paragraph{Input normalisation: the agent-centred frame.}
\index{agent-centred frame}
The network should not learn that trajectories near $x = 40$~m behave
differently from those near $x = -10$~m, nor that a drone heading north
turns differently from one heading east; positions are arbitrary and the
physics is invariant to translation and rotation. The inputs are therefore
\emph{displacements} $\Delta\pos_t = \pos_t - \pos_{t-1}$ expressed in a frame
whose origin is the last observed position and whose $x$-axis points along
the last observed heading (estimated from the last two displacements), and
divided by a typical step length $s$ so that they are of order one:
\begin{equation}
  \vect{x}_t = \frac{1}{s}\,\mat{R}(-\varphi_T)\,(\pos_t - \pos_{t-1}),
  \qquad t = T - T_{\mathrm{obs}} + 2, \dots, T,
  \label{eq:ch20-frame}
\end{equation}
where $\mat{R}(\cdot)$ is a rotation matrix and $\varphi_T$ the last heading.
The targets, the future displacements, are transformed the same way, and
the predicted displacements are summed, rescaled and rotated back to give
world positions. Constant velocity is then a fixed, simple pattern: the
input displacements are all $\approx (\bar{x}, 0)$ and the answer is to
repeat them. The book's implementation makes this explicit by adding the
mean observed displacement $\bar{\vect{x}}$ to every decoder output, so that a
decoder that outputs zero reproduces the constant-velocity prediction and
the network only has to learn the \emph{deviation} from it. In
\cref{sec:ch20-example} the same network trained on absolute positions is
worse than the constant-velocity baseline.

\paragraph{Rolling the decoder out.}
The decoder's input at step $k$ is a displacement. Three choices exist for
where it comes from, and they define three modes of the same network. In
\textbf{free-running} mode the input is the mean $\vect{\mu}_{k-1}$ that the
decoder itself produced one step earlier; this is how the predictor is used
at run time. In \textbf{teacher forcing}\index{teacher forcing} mode, used
during training, the input is the \emph{true} displacement $\vect{y}_{k-1}$:
the decoder is told the correct answer for the previous step and only has
to predict the next one. Teacher forcing makes the early phase of training
fast and stable, because errors do not accumulate along the rollout. Its
weakness is \textbf{exposure bias}\index{exposure bias}: the decoder is
trained on perfect inputs and evaluated on its own imperfect ones, and it
can learn the shortcut of copying its input, which works under teacher
forcing and fails in free-running mode. The experiment of
\cref{sec:ch20-example} shows the effect. The remedies are to train in
free-running mode with the gradient flowing through the feedback, to mix
the two modes with a schedule~\cite{bengio2015scheduled}, or to feed the
decoder no displacement at all and let its state carry everything. In
\textbf{sampling} mode the input is a sample $\vect{\mu}_{k-1} +
\vect{\sigma}_{k-1} \odot \vect{\eps}$ from the predicted Gaussian, which
produces one plausible future per rollout.

\begin{algorithm}[htb]
\caption{Sequence-to-sequence prediction with an encoder and a decoder LSTM}
\label{alg:ch20-seq2seq}
\KwIn{normalised observed displacements $\vect{x}_1, \dots, \vect{x}_{T_{\mathrm{obs}}-1}$, horizon $H$, mode $\in \set{\text{free-running}, \text{teacher forcing}, \text{sampling}}$, true future $\vect{y}_{1:H}$ (teacher forcing only)}
\KwOut{means $\vect{\mu}_{1:H}$ and standard deviations $\vect{\sigma}_{1:H}$ of the future displacements}
\Fn{\PredPredict{$\vect{x}_{1:T_{\mathrm{obs}}-1}$, $H$, mode}}{
  $\vect{h}, \vect{c} \gets \vect{0}$; \quad $\bar{\vect{x}} \gets$ mean of the $\vect{x}_t$\tcp*{the constant-velocity displacement}
  \For{$t = 1, \dots, T_{\mathrm{obs}} - 1$}{$(\vect{h}, \vect{c}) \gets \PredCell_{\mathrm{enc}}(\vect{x}_t, \vect{h}, \vect{c})$\label{alg:ch20-seq2seq:encode}\;}
  $\vect{u} \gets \vect{x}_{T_{\mathrm{obs}}-1}$\tcp*{first decoder input: the last observed displacement}
  \For{$k = 1, \dots, H$}{
    $(\vect{h}, \vect{c}) \gets \PredCell_{\mathrm{dec}}(\vect{u}, \vect{h}, \vect{c})$\label{alg:ch20-seq2seq:decode}\;
    $(\vect{\mu}_k, \log\vect{\sigma}_k) \gets \mat{W}_{\mathrm{out}}\vect{h} + \vect{b}_{\mathrm{out}}$; \quad $\vect{\mu}_k \gets \vect{\mu}_k + \bar{\vect{x}}$\label{alg:ch20-seq2seq:readout}\tcp*{residual over constant velocity}
    \uIf{mode = teacher forcing}{$\vect{u} \gets \vect{y}_k$\;}
    \uElseIf{mode = sampling}{$\vect{u} \gets \vect{\mu}_k + \vect{\sigma}_k \odot \vect{\eps}$, $\vect{\eps} \sim \Normal(\vect{0}, \mat{I})$\;}
    \Else{$\vect{u} \gets \vect{\mu}_k$\label{alg:ch20-seq2seq:feedback}\;}
  }
  \KwRet $(\vect{\mu}_{1:H}, \vect{\sigma}_{1:H})$; positions are the cumulative sums of the means, mapped back to the world frame\;
}
\end{algorithm}

\paragraph{The loss: squared error or Gaussian likelihood.}
\index{loss function!mean squared error}\index{loss function!Gaussian negative log-likelihood}
With a read-out of two numbers per step the natural loss is the
\textbf{mean squared error} of the predicted positions,
\begin{equation}
  \mathcal{L}_{\mathrm{MSE}} = \frac{1}{H}\sum_{h=1}^{H} \norm{\hat{\pos}_{T+h} - \pos_{T+h}}^2,
  \label{eq:ch20-mse}
\end{equation}
averaged over the training batch. It trains a predictor of the
\emph{conditional mean} (\cref{thm:ch20-conditional-mean}) and says nothing
about uncertainty. With a read-out of four numbers per step, a mean
$\vect{\mu}_h$ and a log standard deviation $\log\vect{\sigma}_h$ of the next
displacement, the loss becomes the \textbf{Gaussian negative
log-likelihood} of the true displacements $\vect{y}_h$,
\begin{equation}
  \mathcal{L}_{\mathrm{NLL}} = \frac{1}{H}\sum_{h=1}^{H}\sum_{d \in \set{x, y}}
  \left[ \log\sigma_{h,d} + \frac{(y_{h,d} - \mu_{h,d})^2}{2\sigma_{h,d}^2} \right] + \text{const},
  \label{eq:ch20-nll}
\end{equation}
which is minimised by the conditional mean \emph{and} the conditional
standard deviation: the network learns to say ``I do not know'' where the
training data are unpredictable, for example right after the onset of an
evasive manoeuvre, by predicting a large $\vect{\sigma}_h$. Predicting
$\log\vect{\sigma}$ rather than $\vect{\sigma}$ keeps it positive without a
constraint. A full covariance per step (three numbers, with a correlation)
is a small extension; the original Social LSTM used exactly such a
bivariate Gaussian output~\cite{alahi2016social}. The position uncertainty
at horizon $h$ is then obtained by summing the displacement covariances up
to $h$, which assumes independent step errors and, as
\cref{sec:ch20-example} shows, underestimates the true spread; a correction
factor estimated on the validation set repairs this.

\paragraph{Training.}
\index{Adam optimiser}\index{early stopping}
\Cref{alg:ch20-train} is the loop. The gradient of the loss with respect
to every weight is computed by BPTT through the decoder, through the
hand-over of the state, and through the encoder; in free-running mode it
also flows through the fed-back means (line~\ref{alg:ch20-seq2seq:feedback}
of \cref{alg:ch20-seq2seq}). The weights are updated with
Adam~\cite{kingma2015adam}, which rescales every parameter's step by a
running estimate of its gradient's magnitude and is the default optimiser
for sequence models; the learning rate decays geometrically over the run.
The gradient is clipped to a maximum norm, because a single
badly-conditioned batch can otherwise throw the weights far away.
\textbf{Early stopping} evaluates, after every epoch, the free-running ADE
on the validation set, keeps the weights of the best epoch and stops when
no improvement has been seen for a fixed number of epochs: it is the one
regulariser that every trajectory predictor should have, and it is where
the validation split earns its keep.

\begin{algorithm}[htb]
\caption{Training the predictor with Adam and early stopping}
\label{alg:ch20-train}
\KwIn{training and validation sets of windows, epochs $E$, batch size $B$, learning rates $\eta_0 \to \eta_E$, patience $p$}
\KwOut{the weights $\theta$ of the epoch with the best validation ADE}
\Fn{\PredTrain{train, val}}{
  initialise $\theta$ (forget-gate biases $1$, small read-out); $\theta^\star \gets \theta$; $\mathrm{best} \gets \infty$\;
  \For{epoch $= 1, \dots, E$}{
    $\eta \gets$ geometric interpolation between $\eta_0$ and $\eta_E$\;
    \ForEach{mini-batch of $B$ windows, in random order}{
      normalise to the agent frame (\cref{eq:ch20-frame}); run \PredPredict in free-running (or teacher-forcing) mode\;
      $\mathcal{L} \gets$ \cref{eq:ch20-nll} or \cref{eq:ch20-mse}; $\nabla_\theta\mathcal{L} \gets$ BPTT\label{alg:ch20-train:bptt}\;
      clip $\nabla_\theta\mathcal{L}$ to norm $\le 5$; Adam step with rate $\eta$\label{alg:ch20-train:adam}\;
    }
    $a \gets$ free-running ADE on val\label{alg:ch20-train:val}\;
    \lIf{$a < \mathrm{best}$}{$\mathrm{best} \gets a$; $\theta^\star \gets \theta$; $\mathrm{wait} \gets 0$}
    \lElse{$\mathrm{wait} \gets \mathrm{wait} + 1$; \lIf{$\mathrm{wait} \ge p$}{\KwBreak}}
  }
  \KwRet $\theta^\star$\;
}
\end{algorithm}

\subsection{Interaction: the Social LSTM}
\label{sec:ch20-social}

A drone in a swarm, like a pedestrian in a crowd, moves according to its
neighbours. Alahi et al.~\cite{alahi2016social} made the LSTM predictor
\textbf{interaction-aware} with an idea called \textbf{social
pooling}\index{social pooling}\index{Social LSTM}, shown in
\cref{fig:ch20-social-pooling}. Every agent $i$ has its own LSTM with hidden
state $\vect{h}_i$. Around agent $i$ lay a grid of $N_o \times N_o$ cells of
side $2\rho/N_o$; for every neighbour $j$ within the grid, add its hidden
state $\vect{h}_j$ to the cell in which it lies:
\begin{equation}
  \mat{H}_i[m, n, :] = \sum_{j \ne i} \mathbb{1}\bigl[\pos_j - \pos_i \in \text{cell}(m, n)\bigr]\,\vect{h}_j .
  \label{eq:ch20-pooling}
\end{equation}
The tensor $\mat{H}_i \in \R^{N_o \times N_o \times D}$ is flattened, passed
through a small embedding layer and concatenated with the embedding of the
agent's own displacement as the input of its LSTM cell. Because every
$\vect{h}_j$ summarises the history of agent $j$, agent $i$'s cell sees not
only where its neighbours are but where they are going; because the grid
is centred on $i$, the representation is translation invariant; and because
cells are summed, any number of neighbours fits. The model is trained
jointly on all agents of a scene with a bivariate Gaussian output per step.

\begin{figure}[tb]
  \centering
  \inputfigure{ch20/social-pooling}
  \caption{Social pooling. Left: a grid of cells is centred on agent $i$;
  neighbours inside it ($j_1, j_2, j_3$) contribute their LSTM hidden states,
  a neighbour outside it ($j_4$) is ignored. Middle: the pooling tensor
  $\mat{H}_i$ holds, per cell, the sum of the hidden states of the neighbours
  in that cell. Right: the embedded tensor is fed to agent $i$'s LSTM cell
  together with the embedding of its own displacement.}
  \label{fig:ch20-social-pooling}
\end{figure}

Social pooling is coarse: a neighbour on the boundary between two cells
jumps between them, cells far away are mostly empty, and the grid has a
fixed range. Its successors replaced the grid by a maximum over all
neighbours (Social GAN~\cite{gupta2018social}), by a graph whose edges carry
the relative state of each pair (Trajectron++~\cite{salzmann2020trajectron}),
and, most generally, by attention over the neighbours, the subject of the
next section.

% ---------------------------------------------------------------------
\section{Transformers and attention over agents}
\label{sec:ch20-transformer}

\subsection{Scaled dot-product attention}
\label{sec:ch20-attention}

An LSTM reads the history in order and compresses it into one state
vector; whatever the decoder needs from step $3$ of the history must
survive five updates. \textbf{Attention}\index{attention} removes the bottleneck:
every output looks directly at every input and decides how much each one
matters. Vaswani et al.~\cite{vaswani2017attention} showed that a network
built from attention alone, the \textbf{Transformer}\index{Transformer},
outperforms recurrent networks on sequence transduction, and Giuliari et
al.~\cite{giuliari2021transformer} showed that the same architecture,
applied to one trajectory at a time, beats the LSTM predictors on the
pedestrian benchmarks.

\begin{definition}[Scaled dot-product attention]
\label{def:ch20-attention}
Let $\mat{Q} \in \R^{n_q \times d_k}$ hold $n_q$ \textbf{queries}, $\mat{K} \in
\R^{n_k \times d_k}$ hold $n_k$ \textbf{keys} and $\mat{V} \in \R^{n_k \times d_v}$
the \textbf{values} attached to the keys, one per row. Then
\begin{equation}
  \operatorname{Attention}(\mat{Q}, \mat{K}, \mat{V})
  = \operatorname{softmax}\!\left(\frac{\mat{Q}\mat{K}\T}{\sqrt{d_k}}\right)\mat{V},
  \qquad
  \alpha_{ij} = \frac{\exp(\vect{q}_i\T\vect{k}_j/\sqrt{d_k})}{\sum_{j'}\exp(\vect{q}_i\T\vect{k}_{j'}/\sqrt{d_k})},
  \label{eq:ch20-attention}
\end{equation}
where the softmax acts on every row: the output for query $i$ is
$\sum_j \alpha_{ij}\vect{v}_j$, a convex combination of the values with the
\textbf{attention weights} $\alpha_{ij} \ge 0$, $\sum_j \alpha_{ij} = 1$.
\index{attention!scaled dot-product}
\end{definition}

The dot product $\vect{q}_i\T\vect{k}_j$ scores how relevant input $j$ is to
output $i$; the division by $\sqrt{d_k}$ keeps the scores of order one for
random vectors, so that the softmax does not saturate at initialisation.
\textbf{Multi-head attention}\index{attention!multi-head} runs $M$ such attentions
in parallel on learned projections of the same inputs,
\begin{equation}
  \operatorname{head}_m = \operatorname{Attention}(\mat{X}\mat{W}_m^{Q}, \mat{X}\mat{W}_m^{K}, \mat{X}\mat{W}_m^{V}),
  \qquad
  \operatorname{MultiHead}(\mat{X}) = [\operatorname{head}_1, \dots, \operatorname{head}_M]\,\mat{W}^{O},
  \label{eq:ch20-multihead}
\end{equation}
with $\mat{W}_m^{Q}, \mat{W}_m^{K}, \mat{W}_m^{V} \in \R^{d \times d/M}$ and
$\mat{W}^{O} \in \R^{d \times d}$, so that one head can attend to the last
step, another to the onset of the current manoeuvre, a third to the
nearest neighbour. A Transformer layer adds a position-wise feed-forward
network, residual connections and layer normalisation around the attention
block; the details are in \textcite{vaswani2017attention} and
\textcite{goodfellow2016deep}.

\paragraph{Positional encodings.}
\index{positional encoding}
Attention is a set operation: permuting the inputs permutes the outputs,
and nothing in \cref{eq:ch20-attention} knows that step $7$ comes after
step $6$. The order is injected by adding to the embedding of the input at
step $t$ a \textbf{positional encoding}
\begin{equation}
  \mathrm{PE}(t, 2i) = \sin\!\left(\frac{t}{10000^{2i/d}}\right), \qquad
  \mathrm{PE}(t, 2i + 1) = \cos\!\left(\frac{t}{10000^{2i/d}}\right), \qquad i = 0, \dots, d/2 - 1,
  \label{eq:ch20-pe}
\end{equation}
a vector of sinusoids of geometrically spaced frequencies. Its two useful
properties are that no two steps have the same vector and that the
encoding of step $t + \delta$ is a fixed linear function of the encoding of
step $t$ for every offset $\delta$ (\cref{exr:ch20-pe}), so that a head can
learn to attend ``two steps back'' regardless of $t$. Learned positional
embeddings are the common alternative for short, fixed horizons.

\subsection{An encoder-decoder Transformer for one trajectory}
\label{sec:ch20-tf-single}

Giuliari et al.~\cite{giuliari2021transformer} use the original
encoder-decoder layout. Each observed displacement is embedded by a linear
layer into $\R^d$ and added to its positional encoding; a stack of encoder
layers, each with multi-head self-attention over the $T_{\mathrm{obs}} - 1$
tokens, produces one context vector per observed step. The decoder produces
the future autoregressively: at step $k$ it embeds the displacements
predicted so far, applies \emph{masked} self-attention (a token may attend
only to earlier tokens, so that training can process all $H$ steps in one
pass while remaining causal), then \emph{cross-attention} whose queries come
from the decoder and whose keys and values are the encoder outputs, and
finally a read-out to the next displacement or to the parameters of its
distribution. Training uses teacher forcing on the decoder inputs and the
same losses as \cref{sec:ch20-seq2seq}; prediction is free-running. The
model has no recurrence, so every observed step is one attention hop away
from every predicted step.

\subsection{Attention across agents}
\label{sec:ch20-agents}

Nothing in \cref{eq:ch20-attention} requires the keys to be time steps of
the \emph{same} agent. Let the tokens be the encoded histories of all $N$
agents in the scene (one summary token per agent, or one token per agent
and time step). The query of agent $i$ scored against the keys of all
agents gives weights $\alpha_{ij}$ that say how much agent $j$ matters for
predicting agent $i$, and the output $\sum_j \alpha_{ij}\vect{v}_j$ is a
learned, weighted summary of the neighbourhood: social pooling with a
learned, continuous notion of ``nearby''. Models of this family
(the graph attention of Trajectron++~\cite{salzmann2020trajectron} and the
joint agent-time attention of later Transformers surveyed
in~\cite{rudenko2020human}) are the state of the art for interaction-aware
prediction.

\Cref{fig:ch20-attention} illustrates the mechanism on four drones. To keep
the numbers reproducible without training a Transformer, the query and key
maps are set by hand in \code{toy\_agent\_attention}: with the augmented
features $\vect{k}_j = (\vect{x}_j, \norm{\vect{x}_j}^2, 1)$ and $\vect{q}_i =
(\vect{x}_i, -\tfrac12, -\tfrac12\norm{\vect{x}_i}^2) \cdot \sqrt{d_k}/\ell^2$,
where $\vect{x} = \pos + \tau\vel$ is the position extrapolated $\tau = 2$~s
ahead, the scaled dot product equals $-\norm{\vect{x}_i - \vect{x}_j}^2 / (2\ell^2)$,
so the softmax weights are large for agents whose extrapolated positions
are close. Drone~A splits its attention almost equally between itself and
drone~B, which is on a collision course ($\alpha_{AB} = 0.48$), and gives
drone~C, which is near but diverging, $0.003$ and the distant drone~D
nothing. A trained model learns such maps from data instead of having them
written down, and it learns them per head: one head may score closeness,
another whether the neighbour is ahead or behind, a third its speed.

\begin{figure}[tb]
  \centering
  \inputfigure{ch20/attention}
  \caption{Attention across agents. Left: four drones with their velocities
  (arrows of $2$~s of flight) and their positions extrapolated $2$~s ahead
  (crosses); A and B are on a collision course, C passes nearby, D is far
  away. Right: the attention weights $\alpha_{ij}$ of \cref{eq:ch20-attention}
  with hand-set query and key maps that score the distance between the
  extrapolated positions (rows sum to one). Each drone attends to itself
  and to the drones it will meet; D attends to nobody but itself.}
  \label{fig:ch20-attention}
\end{figure}

\subsection{What Transformers buy and what they cost}
\label{sec:ch20-tf-cost}

Attention helps at long horizons because the decoder can consult any
observed step directly instead of relying on a compressed state, and it
helps with many agents because the same mechanism ranks neighbours by
learned relevance rather than by a grid. Both advantages come with bills.
The first is \emph{data}: a Transformer has more parameters and fewer
built-in assumptions than an LSTM and needs more trajectories to fit them;
on a few hundred recorded manoeuvres a small LSTM, or the constant-velocity
model, is usually the better choice. The second is \emph{compute}: the
scores of \cref{eq:ch20-attention} cost $\bigO{n^2 d}$ for $n$ tokens
(\cref{thm:ch20-complexity}), so a scene with $N$ agents over $T$ steps,
tokenised jointly, costs $\bigO{N^2 T^2 d}$ per layer, against the
$\bigO{N T d^2}$ of one LSTM per agent. For a swarm planner that predicts
every $0.1$~s this matters. The third is \emph{engineering}: the LSTM of
this chapter is $300$ lines of NumPy including the backward pass; a
Transformer wants a framework, and \cref{sec:ch20-implementation} points to
PyTorch for it.

% ---------------------------------------------------------------------
\section{Multimodality and uncertainty}
\label{sec:ch20-multimodal}

\Cref{thm:ch20-conditional-mean} in \cref{sec:ch20-properties} states the
problem precisely: a predictor trained with the squared error converges
to the conditional mean of the future, and when the future has two
distinct modes, left and right around an obstacle, the conditional mean
runs through the obstacle. The Gaussian output of \cref{eq:ch20-nll} does
not fix this; it reports the disagreement as a large $\vect{\sigma}$, which
is honest but still unimodal. Three families of remedies exist.

\paragraph{Mixture density outputs.}
\index{mixture density network}
The read-out can produce, per step, the parameters of a \textbf{mixture of
Gaussians}~\cite{bishop1994mixture}: $C$ weights $\pi_c$ (a softmax over $C$
logits), means $\vect{\mu}_c$ and covariances $\mat{\Sigma}_c$, with the loss
\begin{equation}
  \mathcal{L}_{\mathrm{MDN}} = -\log \sum_{c=1}^{C} \pi_c\, \Normal\bigl(\vect{y}; \vect{\mu}_c, \mat{\Sigma}_c\bigr),
  \label{eq:ch20-mdn}
\end{equation}
computed with the log-sum-exp trick for stability. At a fork the network
learns two components with weights $\approx \tfrac12$; on straight flight
one component takes all the weight. Mixture outputs are the cheapest form
of multimodality and are used by Trajectron++ and many Transformer
predictors.

\paragraph{Sampling-based models.}
\index{generative model!for prediction}
A second approach feeds the decoder a random latent vector $\vect{z}$ and
trains the network so that different $\vect{z}$ produce different plausible
futures. Social GAN~\cite{gupta2018social} does this adversarially: a
discriminator learns to tell generated futures from real ones, and a
\emph{variety loss} that back-propagates only the best of $K$ samples pushes
the generator to cover the modes. Trajectron++~\cite{salzmann2020trajectron}
uses a conditional variational autoencoder with a \emph{discrete} latent
variable, so that each value of $\vect{z}$ is an interpretable mode, and
integrates the sampled controls through a dynamics model so that the
samples are physically feasible. Both are scored with $\mathrm{minADE}_K$;
both need a planner that can consume a set of samples. The autoregressive
sampling mode of \cref{alg:ch20-seq2seq} is the poor man's version: it
produces diverse rollouts but only from a unimodal per-step Gaussian.

\paragraph{Ensembles.}
\index{ensemble}
Training $M$ copies of the predictor with different seeds and averaging
their outputs improves accuracy, and the spread of their predictions is a
usable estimate of the \emph{model} uncertainty, which the per-step
$\vect{\sigma}$ cannot see~\cite{lakshminarayanan2017ensembles}. The cost is
$M$ trainings and $M$ forward passes; five is a common choice.

\paragraph{From a distribution to an occupancy region.}
\index{occupancy region!time-indexed}
Whatever produces the distribution, the planner needs it as geometry. For
a Gaussian prediction at step $h$ with mean $\hat{\pos}_{T+h}$ and covariance
$\mat{\Sigma}_h$, the set of positions the intruder occupies with probability
$p$ is the covariance ellipse of \cref{ch:ch02},
\begin{equation}
  \mathcal{E}_h(p) = \set{\vect{x} : (\vect{x} - \hat{\pos}_{T+h})\T\mat{\Sigma}_h\inv(\vect{x} - \hat{\pos}_{T+h}) \le k_p^2},
  \qquad k_p = \sqrt{-2\ln(1 - p)},
  \label{eq:ch20-ellipse}
\end{equation}
where $k_p$ follows from the $\chi^2$ distribution with two degrees of
freedom ($k_{0.95} = 2.448$, $k_{0.99} = 3.035$). The sequence
$\mathcal{E}_1(p), \dots, \mathcal{E}_H(p)$ is the \textbf{time-indexed
occupancy region} of the intruder. For a mixture, take the union of the
ellipses of the components whose weight exceeds a threshold; for $K$
samples, take the smallest discs, or the grid cells, that contain a
fraction $p$ of the samples at each step. Before any of this is trusted,
its \textbf{calibration}\index{calibration} must be checked on held-out data:
the fraction of true positions that fall inside $\mathcal{E}_h(p)$ should be
about $p$ at every $h$. Learned predictors are often over-confident, and the
experiment below quantifies by how much; the repair is a per-horizon
inflation factor fitted on the validation set. \Cref{sec:ch20-drone} shows
what \orca, \dstarlite and MPC do with the region.
